\section{Method: Training-Adaptive Data Selection with a Reliability Gate (\textsc{TAG})}
\label{sec:method}
\textsc{TAG} separates two questions that are otherwise conflated by
dynamic selection: whether an instruction--response pair provides
credible supervision, and how useful that pair is to the model at its
current training state. A static reliability gate answers the first;
a dynamically refreshed reward and a state-conditioned
representation-anchor term answer the second.
Section~\ref{sec:problem-setup} fixes notation and the selection
rule; Section~\ref{sec:reliability-gate} constructs the gate;
Sections~\ref{sec:method-motivation}--\ref{sec:local-anchor-stability}
develop the two dynamic views and analyze the stability of the
recomputed anchor; Section~\ref{sec:tag-score} interprets the fused
score; Section~\ref{sec:algorithm-implementation} summarizes the
complete algorithm.

\subsection{Problem setup and notation}
\label{sec:problem-setup}

Let \(\mathcal D=\{x_i\}_{i=1}^N\) be a pool, where
\(x_i=(u_i,y_i)\) contains instruction \(u_i\) and response \(y_i\). Let
\(p_{\theta}\) denote the base language model. At refresh step \(t\),
\textsc{TAG} assigns candidate \(x_i\) the score
\begin{equation}
  s_i^{(t)}
  =
  \underbrace{G_i}_{\text{reliability gate}}
  \;\mathcal R_i^{(t)}
  \left(1+\lambda\,
  \widetilde a_i^{(t)}\right),
  \qquad G_i\in[0,1],
  \label{eq:tag-score-method}
\end{equation}
where \(G_i\) is computed once with the base checkpoint \(\theta_0\),
\(\mathcal R_i^{(t)}\) is a dynamically refreshed
difficulty--uncertainty reward,
\(\widetilde a_i^{(t)}\in[0,1]\) is the candidate's normalized
projection onto the current representation anchor, and
\(\lambda\ge 0\) is a fixed anchor-modulation strength, confining the
anchor factor to \([1,1+\lambda]\). Given budget \(B\le N\),
the implementation uses the lexicographic key
\begin{equation}
  \kappa_i^{(t)}
  :=
  \begin{cases}
    2+\operatorname{rank}_{01}\!\left(s_i^{(t)}\right), & G_i>0,\\
    \operatorname{rank}_{01}\!\left(\widehat{\Delta}_i\right), & G_i=0,
  \end{cases}
  \qquad
  \mathcal S_t
  :=\operatorname{TopB}_{x_i\in\mathcal D}\kappa_i^{(t)},
  \label{eq:tag-selection-key}
\end{equation}
where \(\operatorname{rank}_{01}\) is the pool-wise midrank normalized
to \([0,1]\), and \(\widehat{\Delta}_i\) is the gate evidence defined
in Section~\ref{sec:reliability-gate}. Thus every candidate with a
positive gate precedes the exact-zero block. If that admissible set
fills the budget, Eq.~\eqref{eq:tag-selection-key} reduces to top-\(B\)
selection by \(s_i^{(t)}\); otherwise, it backfills the remaining slots
with the least unreliable zero-gate candidates. The gate is fixed
across refreshes, whereas usefulness is re-evaluated as
\(\theta_{t-1}\) evolves.
\paragraph{Notation.}
We write \(\mathcal L_i^{(t)}\) for response-token cross-entropy,
\(\mathcal H_i^{(t)}\) for mean predictive entropy, and
\(w^{(t)}\) for their variance-ratio coefficient. Superscript \(0\)
denotes a quantity evaluated once at \(\theta_0\); superscript \(t\)
denotes a quantity refreshed under \(\theta_{t-1}\).

\subsection{Reliability gate}
\label{sec:reliability-gate}

\paragraph{Design requirements.}
The gate must resolve the ambiguity that defeats dynamic utility
signals in low-quality pools: loss and entropy rise both for
genuinely hard clean pairs and for corrupted ones, so a
difficulty-driven selector cannot separate ``hard because
informative'' from ``hard because broken''
(Section~\ref{sec:intro}). Fluency is no remedy---a mismatched or
subtly wrong response is often perfectly fluent, which is precisely
why perplexity-style filters pass it. We therefore require a signal
that (i) measures pair-level credibility on an axis orthogonal to
model-state utility, (ii) uses no external judge, reward model, or
per-pair trusted reference response---the only supervision assumed
is a small pool of responses presumed clean, used once to calibrate
the null---and (iii) is cheap enough to run once over the full pool
and cache. The construction below meets these
requirements with two forward passes per pair at the base
checkpoint.

\paragraph{Instruction-side counterfactual contrast.}
What corruption destroys is the predictive dependence
\(u_i\!\to\!y_i\). Intervening on the response side cannot isolate
this dependence: a mismatched \(y_i\) remains a fluent, valid answer
to \emph{some} instruction, and the trusted replacement for a
suspect \(y_i\) is exactly the supervision that is unknown. The
instruction side, by contrast, offers a valid negative control for
free. We hold \(y_i\) fixed and replace \(u_i\): for each
\(x_i=(u_i,y_i)\), a derangement within the same response-length
stratum assigns a counterfactual instruction \(u_i^-\), so that no
original pairing is retained and the negative condition is a
\emph{real} instruction of comparable scale. The counterfactual is
used only for scoring and never enters training. With
\(u_i^+=u_i\) and \(K_i\) the response-token support shared by the
two prompts, define at the frozen base checkpoint
\begin{equation}
\begin{aligned}
  L_i^{\pm}
  &:=-\sum_{k=1}^{K_i}
  \log p_{\theta_0}(y_{ik}\mid u_i^{\pm},y_{i,<k}),
  \\
  \overline{\Delta}_i
  &:=1-\frac{L_i^+}{L_i^-}.
\end{aligned}
  \label{eq:gate-counterfactual-support}
\end{equation}
A positive \(\overline{\Delta}_i\) says that the paired instruction
predicts the fixed response better than a matched in-pool
alternative. Because the reference condition is a real instruction
rather than an empty or generic one, loss mass that does not depend
on the condition---fluency, boilerplate, length effects---enters
both terms equally and is differenced out of the contrast; its
residual dilutive effect on long responses is removed by the
calibration below. This matched negative control distinguishes the contrast
from conditional-versus-unconditional difficulty scores such as
IFD~\cite{q2q}, while following the broader principle of measuring
model-relative input contribution through changes in conditional
likelihood~\cite{ethayarajh2022usable}. This is an
intervention-style diagnostic of conditional support, not a claim
of causal identification, and it measures
instruction--response dependence rather than certifying factual
correctness.

\paragraph{From evidence to a calibrated, zero-anchored gate.}
Two refinements turn the raw contrast into a reusable weight; each
targets a specific failure mode of the plain sequence average.
\emph{(a)~Conservative span evidence.} A sequence-level mean can
hide a locally unsupported segment---a wrong final answer or a
truncated tail in an otherwise supported response. We therefore
repeat the contrast in Eq.~\eqref{eq:gate-counterfactual-support}
over short contiguous response spans and let \(q_i\) be the minimum
of the sequence-level and admissible span-level contrasts, asking
whether the instruction supports the response \emph{throughout}
while excluding short or uninformative spans.
\emph{(b)~Length-conditioned calibration.} A minimum over more
spans has a lower null value even for clean responses, so an
uncalibrated threshold would systematically penalize long
responses. We calibrate \(q_i\) against clean responses with a
comparable span count \(M_i\), whose lower envelope we denote
\(\mu(M_i)\), center the evidence, and map it to the gate:
\begin{equation}
\begin{aligned}
  \widehat{\Delta}_i&:=q_i-\mu(M_i),\\
  G_i
  &:=c_i\left[2\sigma\!\left(\widehat{\Delta}_i/s_G\right)-1\right]_+
  \in[0,1].
\end{aligned}
  \label{eq:reliability-gate}
\end{equation}
Here \(s_G\) is fixed from the same reference distribution and
\(c_i\in[0,1]\) attenuates visibly incomplete responses. Exact span
construction, calibration, completeness checks, and edge-case
handling are given in Appendix~\ref{app:gate-implementation}.

\paragraph{Properties and role.}
The mapping itself is deliberately simple; its value lies in three
properties that the fused rule in Eq.~\eqref{eq:tag-score-method}
exploits. First, the gate is \emph{zero-anchored}: evidence at or
below the clean envelope maps to exactly \(G_i=0\), so an
unsupported pair is vetoed multiplicatively rather than merely
down-weighted, and no amount of dynamic utility can compensate
(Section~\ref{sec:tag-score}). Second, it is \emph{calibrated}:
centering against the length-conditioned clean envelope makes
\(G_i\) comparable across response lengths, which is what allows
one continuous weight to gate a heterogeneous pool. Third, it is
\emph{static and cached}: computed once under \(\theta_0\), it adds
no per-refresh cost, so reliability screening scales independently
of how often the dynamic views are refreshed. In combination, the
matched in-pool counterfactual, the conservative span minimum, and
the zero-anchored length-conditioned calibration form a
pre-selection reliability stage that existing dynamic selectors
lack (Section~\ref{sec:related}); none of the three components
requires per-pair labels, external judges, or selector training
beyond the one-time clean-reference calibration above. The gate decides
admissibility only; the next two views measure usefulness under the
current checkpoint.

\subsection{Difficulty--uncertainty carrier and its structural
limitation}
\label{sec:method-motivation}

\paragraph{Difficulty--uncertainty carrier.}
At refresh step \(t\), \(\mathcal{L}_i^{(t)}\) and
\(\mathcal{H}_i^{(t)}\) are computed under the current checkpoint
\(\theta_{t-1}\). Following dynamic difficulty--uncertainty selection,
we define the carrier reward
\begin{equation}
  \mathcal{R}_i^{(t)}
  =
  w^{(t)}\mathcal{L}_i^{(t)}
  +
  \bigl(1-w^{(t)}\bigr)\mathcal{H}_i^{(t)},
  \label{eq:composite}
\end{equation}
where the variance-ratio weight is
\begin{equation}
  w^{(t)}
  =
  \frac{\operatorname{Var}_{\mathcal D}(\mathcal L^{(t)})}
       {\operatorname{Var}_{\mathcal D}(\mathcal L^{(t)})
       +
        \operatorname{Var}_{\mathcal D}(\mathcal H^{(t)})
       +
        \varepsilon}.
  \label{eq:w}
\end{equation}
In \textsc{TAG}, this loss--entropy reward is kept as a deterministic
carrier modulated by the representation anchor~\cite{kung2023activeit,dataagent}.

\paragraph{Structural limitation.}
The variance-ratio carrier can mix within-family informativeness with
between-family scale effects in heterogeneous pools. To see this,
suppose for analysis only that \(\mathcal D\) is partitioned into
task-family groups \(\{\mathcal D_g\}_{g=1}^G\), and let
\(\mu_g^{(t)}=\mathbb E_{x\in\mathcal D_g}[\mathcal L^{(t)}(x)]\).
By the law of total variance,
\begin{equation}
  \operatorname{Var}_{\mathcal D}(\mathcal L^{(t)})
  =
  \underbrace{
  \mathbb E_g
  \!\left[
  \operatorname{Var}_{\mathcal D_g}(\mathcal L^{(t)})
  \right]
  }_{\text{\scriptsize within-family}}
  +
  \underbrace{
  \operatorname{Var}_g(\mu_g^{(t)})
  }_{\text{\scriptsize between-family}} .
  \label{eq:total-variance}
\end{equation}
The same decomposition holds for \(\mathcal H^{(t)}\). Thus, in
multi-task pools, \(w^{(t)}\) may reflect task-family scale differences
in addition to per-sample difficulty and uncertainty. \textsc{TAG}
therefore keeps \(\mathcal R_i^{(t)}\) as the carrier but modulates
it with candidate projections onto a representation anchor recomputed
from the current model state.


\subsection{State-conditioned representation anchor and candidate
projection}
\label{sec:representation-anchor}

\paragraph{Principal-direction construction.}
Let \(h_l^{(k)}(x;\theta)\in\mathbb R^d\) denote the hidden state of
sequence \(x\) at layer \(l\) and token position \(k\). We define the
within-sequence representation displacement and sequence-mean
activation as
\begin{equation}
  \Delta h_l(x;\theta)
  :=
  h_l^{(K_x)}(x;\theta)-h_l^{(1)}(x;\theta),
  \label{eq:delta}
\end{equation}
and
\begin{equation}
  \bar h_l(x;\theta)
  :=
  \frac{1}{K_x}\sum_{k=1}^{K_x} h_l^{(k)}(x;\theta).
  \label{eq:meanact}
\end{equation}

At refresh step \(t\), we sample a probe subset
\(\widetilde{\mathcal D}_t\subset\mathcal D\). For each retained layer
\(l\), we compute the mean probe delta
\begin{equation}
  \overline{\Delta h}_l^{(t)}
  :=
  \frac{1}{|\widetilde{\mathcal D}_t|}
  \sum_{x\in\widetilde{\mathcal D}_t}
  \Delta h_l(x;\theta_{t-1}),
  \label{eq:anchor-delta-mean}
\end{equation}
and the centred empirical covariance
\begin{equation}
\begin{aligned}
  \widehat{\Sigma}^{(t)}_l
  :=
  \frac{1}{|\widetilde{\mathcal D}_t|}
  \sum_{x\in\widetilde{\mathcal D}_t}
  &\bigl(\Delta h_l(x;\theta_{t-1})
  -\overline{\Delta h}_l^{(t)}\bigr) \\
  &\bigl(\Delta h_l(x;\theta_{t-1})
  -\overline{\Delta h}_l^{(t)}\bigr)^\top .
\end{aligned}
\label{eq:anchor-covariance}
\end{equation}
The top unit PCA direction is
\begin{equation}
  \tilde v_l^{(t)}
  \in
  \arg\max_{\|v\|=1}
  v^\top \widehat{\Sigma}^{(t)}_l v .
  \label{eq:anchor-eigenvector}
\end{equation}
By the Rayleigh--Ritz characterization,
\(\tilde v_l^{(t)}\) is the axis that maximizes empirical variance
among the centered probe displacements. It is not assumed to be a
semantic capability direction. Because PCA directions are
sign-indeterminate, we orient \(\tilde v_l^{(1)}\) toward
\(\overline{\Delta h}_l^{(1)}\) at the first refresh and, for
\(t>1\), toward \(v_l^{(t-1)}\). We use the resulting oriented
principal representation-displacement direction \(v_l^{(t)}\) as a
state-conditioned representation anchor.

\paragraph{Candidate anchor projection.}
For candidate \(x_i\), the raw anchor-projection score aggregates
layer-wise signed projections of the candidate's sequence-mean
activation onto the current anchor:
\begin{equation}
  a_i^{(t)}
  =
  \sum_{l=1}^{L}
  \left\langle
  \bar h_l(x_i;\theta_{t-1}),
  v_l^{(t)}
  \right\rangle .
  \label{eq:raw-anchor-projection}
\end{equation}
Although \(v_l^{(t)}\) has unit norm, \(\bar h_l\) is not
normalized. Thus \(a_i^{(t)}\) is a signed projection, not a cosine
similarity; it retains both directional and activation-magnitude
information. The raw scores are min--max normalized over the candidate
pool. Let
\(m_t=\min_{j\in\mathcal D}a_j^{(t)}\) and
\(M_t=\max_{j\in\mathcal D}a_j^{(t)}\). We set
\begin{equation}
  \widetilde a_i^{(t)}
  =
  \begin{cases}
  \dfrac{a_i^{(t)}-m_t}{M_t-m_t},
  & M_t-m_t\ge 10^{-8},\\[4pt]
  1/2, & \text{otherwise}.
  \end{cases}
  \label{eq:anchor-projection-norm}
\end{equation}
The second case makes the representation-anchor modulation uniform,
preserving the carrier-only ranking. Because both the anchor and the
candidate activations are recomputed under \(\theta_{t-1}\), the
modulation tracks the current model state; its usefulness therefore
hinges on whether the recomputed direction is stable between
consecutive refreshes, which we analyze and verify next.


\begin{figure}[t]
\centering
\begin{tikzpicture}[scale=0.95, every node/.style={font=\small}]
  \draw[gray!30, very thin, -{Latex[length=2mm]}] (-0.1,0) -- (6.4,0);
  \draw[gray!30, very thin, -{Latex[length=2mm]}] (0,-0.1) -- (0,3.8);
  \node[gray!50, font=\scriptsize] at (6.5,-0.18) {$\mathbb{R}^d$};

  \filldraw[black] (0,0) circle (1.2pt);

  \foreach \x/\y in
    {3.8/1.2, 4.2/1.6, 3.5/1.8, 4.5/2.0, 3.0/1.3, 4.0/2.3, 3.7/1.5, 4.6/1.4, 3.3/1.6}
  {
    \draw[-{Latex[length=1.6mm]}, orange!60, thin, opacity=0.75]
      (0,0) -- (\x,\y);
  }

  \node[orange!80!black, font=\footnotesize, align=left, anchor=north west] at (3.5,0.85)
    {Probe deltas\\
     $\{\Delta h_l(x;\theta_{t-1})\}_{x\in\widetilde{\mathcal{D}}_t}$};

  \draw[-{Latex[length=2.6mm]}, blue!75!black, ultra thick]
    (0,0) -- (5.0,2.0);
  \node[blue!75!black, font=\bfseries] at (5.35,2.25) {$v_l^{(t)}$};

  \draw[-{Latex[length=2.2mm]}, blue!35, dashed, semithick]
    (0,0) -- (5.0,1.65);
  \node[blue!50, font=\scriptsize] at (5.40,1.45) {$v_l^{(t-1)}$};

  \node[gray!60!black, font=\scriptsize, align=left] at (6.5,1.5)
    {previous\\
     sign-calibrated\\
     anchor};

  \filldraw[red!75!black] (2.2,2.85) circle (2pt);
  \node[red!75!black, font=\footnotesize] at (1.45,3.05)
    {$\bar h_l(x_i;\theta_{t-1})$};

  \draw[red!75!black, dotted, thick] (2.2,2.85) -- (2.88,1.15);
  \filldraw[red!75!black] (2.88,1.15) circle (1.5pt);

  \node[red!75!black, font=\scriptsize, align=left, anchor=west]
    at (0.70,1.45)
    {projection score\\[-1pt]
     $\langle \bar h_l(x_i;\theta_{t-1}), v_l^{(t)}\rangle$};
\end{tikzpicture}

\caption{\textbf{State-conditioned representation anchor at refresh
step~\(t\).}
\textnormal{One-layer schematic. \(\Delta h_l\): probe
representation displacements;
\(v_l^{(t)}\): current anchor; \(v_l^{(t-1)}\): previous anchor;
\(\bar h_l(x_i;\theta_{t-1})\): candidate mean activation. Dotted
line: the layer-wise projection contribution
\(\langle \bar h_l(x_i;\theta_{t-1}), v_l^{(t)}\rangle\), which is
summed in Eq.~\eqref{eq:raw-anchor-projection} to form the raw
anchor-projection score \(a_i^{(t)}\).}}
\label{fig:anchor}
\end{figure}


\subsection{Anchor stability: local analysis and empirical
verification}
\label{sec:local-anchor-stability}

\paragraph{Local stability bound.}
The following Davis--Kahan-style bound records the local sensitivity
of the recomputed PCA anchor. It is used only to motivate the
diagnostic evaluated below; the algorithm itself does not use this
bound.

\begin{theorem}[State-conditioned anchor stability]
\label{thm:anchor-stability}
For notational simplicity, write
\(\Sigma_l^{(t)}=\widehat{\Sigma}_l^{(t)}\), and let \(\eta_{t-1}\)
denote the optimizer learning rate in effect at refresh step \(t\).
Suppose that, for layer \(l\) at refresh step \(t\),
\begin{equation}
  \left\|
  \Sigma_l^{(t)}-\Sigma_l^{(t-1)}
  \right\|_F
  \le
  C_{\Sigma,l}\eta_{t-1},
  \label{eq:cov-drift-assumption}
\end{equation}
and
\begin{equation}
  \lambda_1(\Sigma_l^{(t)})
  -
  \lambda_2(\Sigma_l^{(t-1)})
  \ge
  \gamma_l>0 .
  \label{eq:gap-assumption}
\end{equation}
Under the temporal sign-calibration convention
\(\langle v_l^{(t)},v_l^{(t-1)}\rangle\ge 0\), the anchor displacement
satisfies
\begin{equation}
  \left\|
  v_l^{(t)}-v_l^{(t-1)}
  \right\|_2
  \le
  \frac{2C_{\Sigma,l}}{\gamma_l}\eta_{t-1}.
  \label{eq:anchor-stability-bound}
\end{equation}
\end{theorem}

The proof and the cross-time eigengap convention are deferred to
Appendix~\ref{app:proof-stability}. The bound is conditional: its
constants depend on the realized covariance trajectory, so it is not
a training-time certificate, and it says nothing about candidate
projections or the full ranking. What it does supply is an
observable prediction---anchor displacement at the optimizer
scale---which we now check on a realized trajectory using a
post-hoc verifier that re-extracts each layer's anchor at fixed
intervals during training (protocol below).

\paragraph{Step-to-step anchor displacement.}
\label{sec:empirical-anchor-stability}
For each retained layer \(l\) and consecutive refreshes
\(t-1\!\to\!t\), we measure the step-to-step anchor displacement
\begin{equation}
  d_l^{(t)}
  :=
  \left\|
  v_l^{(t)}-v_l^{(t-1)}
  \right\|_2 ,
  \label{eq:dstep}
\end{equation}
where \(v_l^{(t)}\) denotes the verifier-extracted anchor after
adjacent-refresh sign calibration. Equivalently, if
\(\tilde v_l^{(t)}\) denotes the raw uncalibrated top eigenvector, then
\begin{equation}
  d_l^{(t)}
  =
  \min_{\sigma\in\{\pm1\}}
  \left\|
  \tilde v_l^{(t)}
  -
  \sigma \tilde v_l^{(t-1)}
  \right\|_2
  =
  \sqrt{
  2
  -
  2
  \left|
  \left\langle
  \tilde v_l^{(t)},\tilde v_l^{(t-1)}
  \right\rangle
  \right|
  } .
  \label{eq:dstep-sign-invariant}
\end{equation}
Thus \(d_l^{(t)}\in[0,\sqrt{2}]\). Small values indicate that the
anchor direction is stable across refreshes, while values near
\(\sqrt{2}\) indicate nearly orthogonal consecutive anchors even
after sign calibration. Under the assumptions of
Theorem~\ref{thm:anchor-stability}, the same transition obeys the
bound in Eq.~\eqref{eq:anchor-stability-bound}. We use
\(d_l^{(t)}\) only as a post-hoc diagnostic; it is not used by
Algorithm~\ref{alg:tag} to rank candidates or select data.

\paragraph{Diagnostic protocol.}
We apply a post-hoc verifier to Qwen2.5-0.5B + LoRA training on
Alpaca-GPT4, a small instrumented setting chosen so that per-refresh
anchor extraction over all layers is affordable. Every 25 optimizer
steps, the verifier re-extracts each layer's anchor from a fixed
2,000-sample diagnostic probe and records \(d_l^{(t)}\). This probe
is fixed across refreshes and separate from the resampled
training-time probe \((n_p=1024)\), so \(d_l^{(t)}\)
removes cross-refresh probe-sampling variance and reduces finite-sample
eigenvector noise.

We discard the first 9 refreshes as warm-up and exclude boundary
layers 1 and 24, whose representation displacements are dominated by
embedding-adjacent and output-projection effects. Empirically, these
effects produce near-degenerate local covariance spectra, weakening the
spectral-separation premise of Theorem~\ref{thm:anchor-stability}. The
diagnostic is therefore evaluated on 22 inner decoder layers over 68
post-warmup transitions, giving \(22\times68=1{,}496\) valid
layer--refresh pairs. The global plug-in envelope
\((2\widehat C_{\Sigma}/\gamma_{\min})\eta=46.24\) is too conservative
to be informative at the scale of \(d_l^{(t)}\), as it exceeds the
maximum sign-calibrated displacement \(\sqrt{2}\). The meaningful
comparison is the trajectory-aware Davis--Kahan envelope of
Fig.~\ref{fig:anchor-drift}, which accumulates realized
per-transition drift-to-gap ratios into an upper envelope on the
remaining drift to the final anchor,
\(\|v_l^{(t)}-v_l^{(T)}\|\); its construction is given in
Appendix~\ref{app:anchor-stability-diagnostics}. The measured
\textsc{TAG} drift stays roughly two orders of magnitude below this
envelope throughout training. Unlike
Table~\ref{tab:dstep-spread}, which aggregates only the
\(22\times68\) retained pairs, the figure traces the full
24-layer, 78-point trajectory including warm-up, so the two views
are complementary rather than identical.

\begin{table}[t]
\centering
\caption{Anchor-displacement diagnostic for \textsc{TAG} and matched
CO~\((\lambda{=}0)\). The invariant statistic is
\(d_l^{(t)}\in[0,\sqrt{2}]\). Med. and peak are transition-level layer
summaries aggregated over post-warmup refreshes; worst step is the
maximum retained layer--refresh pair.}
\label{tab:dstep-spread}

\small
\setlength{\tabcolsep}{4pt}
\begin{tabular}{@{}lccccc@{}}
\toprule
 & \multicolumn{2}{c}{\(d^{(t)}\) over layers}
 & worst & sign-flip & \\
\cmidrule(lr){2-3}
Method & med. & peak & step & rate & Status \\
\midrule
\textsc{TAG} \((\lambda=1)\)
 & \(0.034\) & \(0.166\) & \(0.623\) & \(1.7\%\) & \textbf{PASS} \\
CO \((\lambda=0)\)
 & \(0.215\) & \(0.527\) & \(1.405\) & \(\mathbf{11.8\%}\) & \textbf{FAIL} \\
\bottomrule
\end{tabular}
\end{table}

\begin{figure}[t]
\centering
\resizebox{0.97\columnwidth}{!}{%
\begin{tikzpicture}[
  every node/.style={font=\scriptsize},
  >=Latex
]


\path[use as bounding box] (-0.55,-0.95) rectangle (9.55,6.90);

% Y-axis grid lines.
\foreach \logv/\ylab in {-3/$10^{-3}$, -2/$10^{-2}$, -1/$10^{-1}$,
                          0/$10^{0}$,  1/$10^{1}$,   2/$10^{2}$, 3/$10^{3}$} {
  \pgfmathsetmacro{\ycoord}{\logv + 3}
  \draw[gray!30, very thin] (0, \ycoord) -- (9, \ycoord);
  \node[anchor=east, font=\tiny] at (-0.05, \ycoord) {\ylab};
}

% X-axis tick labels.
\foreach \step/\xc in {25/0.00, 500/2.34, 1000/4.68, 1500/6.95, 1950/9.00} {
  \draw[gray!50, very thin] (\xc, 0) -- (\xc, -0.08);
  \node[anchor=north, font=\tiny] at (\xc, -0.10) {\step};
}

% Measured Davis--Kahan cumulative tail bound for TAG.
\draw[blue!70!black, very thick, dashed]
  plot[smooth] coordinates {
    (0.000,5.11) (0.467,5.05) (0.935,4.98) (1.403,4.92) (1.870,4.86)
    (2.338,4.78) (2.805,4.70) (3.273,4.61) (3.740,4.52) (4.208,4.45)
    (4.675,4.37) (5.143,4.28) (5.610,4.20) (6.078,3.97) (6.545,3.81)
    (7.013,3.63) (7.481,3.46) (7.948,3.31) (8.416,3.05) (9.000,2.71)
  };

% Matched CO (TAG, lambda=0) control: same post-hoc diagnostic,
% anchor multiplier disabled during selection.
\draw[gray!55!black, thick, densely dotted]
  plot[smooth] coordinates {
    (0.000,3.146) (0.467,3.114) (0.935,3.097) (1.403,3.079) (1.870,3.130)
    (2.338,3.041) (2.805,3.079) (3.273,3.061) (3.740,3.146) (4.208,3.041)
    (4.675,3.097) (5.143,3.079) (5.610,3.114) (6.078,3.041) (6.545,3.130)
    (7.013,3.114) (7.481,3.146) (7.948,3.130) (8.416,3.161) (9.000,3.146)
  };
\foreach \p in {
  (0.000,3.146), (0.467,3.114), (0.935,3.097), (1.403,3.079), (1.870,3.130),
  (2.338,3.041), (2.805,3.079), (3.273,3.061), (3.740,3.146), (4.208,3.041),
  (4.675,3.097), (5.143,3.079), (5.610,3.114), (6.078,3.041), (6.545,3.130),
  (7.013,3.114), (7.481,3.146), (7.948,3.130), (8.416,3.161), (9.000,3.146)
} {
  \fill[gray!55!black] \p circle (1.3pt);
}

% Circle end-of-schedule CO points for diagnostic contrast.
\foreach \p in {(8.416,3.161), (9.000,3.146)} {
  \draw[gray!55!black, thick] \p circle (4.0pt);
}

% TAG measured cumulative drift to final anchor.
\draw[red!75!black, thick]
  plot[smooth] coordinates {
    (0.000,2.79) (0.467,2.33) (0.935,2.31) (1.403,2.27) (1.870,2.24)
    (2.338,2.14) (2.805,2.12) (3.273,1.98) (3.740,1.88) (4.208,1.95)
    (4.675,1.81) (5.143,1.80) (5.610,1.77) (6.078,1.62) (6.545,1.42)
    (7.013,1.33) (7.481,1.03) (7.948,0.97) (8.416,0.85) (9.000,0.77)
  };
\foreach \p in {
  (0.000,2.79), (0.467,2.33), (0.935,2.31), (1.403,2.27), (1.870,2.24),
  (2.338,2.14), (2.805,2.12), (3.273,1.98), (3.740,1.88), (4.208,1.95),
  (4.675,1.81), (5.143,1.80), (5.610,1.77), (6.078,1.62), (6.545,1.42),
  (7.013,1.33), (7.481,1.03), (7.948,0.97), (8.416,0.85), (9.000,0.77)
} {
  \fill[red!75!black] \p circle (1.5pt);
}

% Axes.
\draw[->, thick] (0,0) -- (9.25,0);
\draw[->, thick] (0,0) -- (0,6.40);

\node at (4.50,-0.55)
  {refresh step $t$ (cosine schedule, 78 points)};

\node[anchor=west, align=left, font=\scriptsize] at (-0.50,6.65)
  {\(\|v_l^{(t)} - v_l^{(T)}\|\) (median over 24 layers, log scale)};

% Legend.
\draw[blue!70!black, very thick, dashed] (4.55,6.30) -- (5.20,6.30);
\node[anchor=west, font=\tiny] at (5.30,6.30)
  {Measured Davis--Kahan tail:
  $\sum_{s \ge t} 2\|\Sigma^{(s+1)}{-}\Sigma^{(s)}\|_F/\gamma^{(s)}$};

\draw[red!75!black, thick] (4.55,5.85) -- (5.20,5.85);
\node[anchor=west, font=\tiny] at (5.30,5.85)
  {\textsc{TAG} measured $\|v^{(t)}{-}v^{(T)}\|$};

\draw[gray!55!black, thick, densely dotted] (4.55,5.40) -- (5.20,5.40);
\node[anchor=west, font=\tiny] at (5.30,5.40)
  {CO ($\lambda{=}0$): same post-hoc diagnostic};

\end{tikzpicture}%
}


\caption{\textbf{Representation-anchor stability diagnostic.}
\textnormal{Median layerwise drift to the final anchor,
$\|v_l^{(t)} - v_l^{(T)}\|$ (log scale; $24$ decoder layers,
$78$ cosine-scheduled refresh points; Qwen2.5-0.5B + LoRA).
\emph{Red}: \textsc{TAG} measured drift.
\emph{Blue dashed}: empirical Davis--Kahan tail
$\sum_{s\ge t} 2\|\Sigma^{(s+1)}-\Sigma^{(s)}\|_F/\gamma^{(s)}$.
\emph{Grey dotted}: matched
CO~($\lambda{=}0$) control;
circled markers denote the final two refreshes.}}
\label{fig:anchor-drift}
\end{figure}

\paragraph{Matched Carrier-Only diagnostic.}
We repeat the verifier on a matched Carrier-Only control,
CO~\((\lambda{=}0)\), which removes the anchor factor from the
selection score. CO and \textsc{TAG} share the same cached gate,
optimizer schedule, fixed diagnostic probe, refresh interval, and verifier
protocol; they differ only in the selection score. This matched design separates
representation-anchor modulation from the optimizer schedule and
diagnostic procedure.

For each refresh transition, we summarize \(d_l^{(t)}\) over the
22 retained layers by the layer-wise median and maximum.
Table~\ref{tab:dstep-spread} reports the median over transitions of
these summaries, denoted ``med.'' and ``peak,'' respectively, together
with the worst step \(\max_{l,t}d_l^{(t)}\) over all 1,496 retained
pairs. The sign-flip rate is the fraction of retained pairs satisfying
\(\langle \tilde v_l^{(t)}, v_l^{(t-1)} \rangle < 0\). This rate is
only an orientation diagnostic; the sign-invariant stability statistic
is \(d_l^{(t)}\). For table status, a run passes if no retained pair
exceeds \(d_l^{(t)}=1\), corresponding to a sign-calibrated angle below
\(60^\circ\). This descriptive criterion is not a theorem-level
guarantee; it flags transitions approaching the near-orthogonal regime.
The contrast in Table~\ref{tab:dstep-spread} is descriptive but
informative: the \textsc{TAG} trajectory stays well inside the stable
regime, while the carrier-only trajectory repeatedly approaches
near-orthogonal anchor transitions. The stability that the bounded
modulation relies on is thus a property the anchored run empirically
maintains, not merely an assumption.


\subsection{Fused score: a zero-anchored product rule}
\label{sec:tag-score}

Equation~\eqref{eq:tag-score-method} fuses the three views
multiplicatively. The product form itself is classical---product-rule
and log-opinion-pool combination of evidence goes back at least to
combining-classifier and expert-aggregation
theory~\cite{kittler1998combining,genest1986combining}---and in log
space the score is simply the additive pool
\(\log s_i^{(t)}=\log G_i+\log\mathcal R_i^{(t)}
+\log\bigl(1+\lambda\,\widetilde a_i^{(t)}\bigr)\).
What the design contributes is the deliberate \emph{asymmetry of the
three roles} within that pool. The gate is absolute and
zero-anchored: \(G_i=0\) annihilates the score and, via the
selection key of Eq.~\eqref{eq:tag-selection-key}, cannot be
re-admitted ahead of any positive-gate candidate, whatever the
dynamic terms---the veto is non-compensatory. Additive or
averaging fusion cannot provide this behavior: a sufficiently large
utility term always outweighs a zero reliability score,
re-admitting exactly the corrupted-but-difficult candidates the
gate exists to remove. The same point separates the design from
trusted multi-view classification, which re-weights or averages
conflicting views~\cite{tmc2021,rcml2024}: in data selection,
unreliability must be able to reach the floor, not merely lose
weight. The two dynamic factors are, by design, weaker: the carrier
\(\mathcal R_i^{(t)}\) is pool-relative and carries the dynamic
ranking, while the anchor factor is confined to
\([1,1+\lambda]\), so it modulates the carrier by at most a factor
of \(1{+}\lambda\) and can neither veto nor dominate. The views
also differ in acquisition schedule---\(G_i\) once at \(\theta_0\),
carrier and anchor at every refresh---so reliability screening adds
no recurring cost. Finally, each view is independently ablatable:
setting \(G_i\equiv 1\) removes the gate without changing the
underlying selector, and setting \(\lambda=0\) removes
representation-anchor modulation, leaving the carrier-only
configuration CO used as the matched control above and as a
baseline in Section~\ref{sec:experiments}.

\subsection{Algorithm and implementation}
\label{sec:algorithm-implementation}

Algorithm~\ref{alg:tag} summarizes the complete procedure. The gate
requires two base-checkpoint forward passes---one under the observed
instruction and one under its counterfactual---and is then cached. Each
refresh uses the current checkpoint only for the reward and anchor.

\begin{algorithm}[t]
  \caption{\textsc{TAG}: reliability-gated training-adaptive selection}
  \label{alg:tag}
  \begin{algorithmic}[1]
  \STATE \textbf{Input:} Pool $\mathcal D$; clean gate reference $\mathcal D_{\mathrm{ref}}$; model $p_{\theta_0}$; budget $B$; anchor strength $\lambda$; refreshes $T$
  \STATE Construct condition-side counterfactuals and compute and cache $G_i$ by Section~\ref{sec:reliability-gate}; see Appendix~\ref{app:gate-implementation}.
  \FOR{$t = 1, \dots, T$}
    \STATE Sample probe $\widetilde{\mathcal D}_t\subset\mathcal D$ uniformly.
    \STATE Under $\theta_{t-1}$, compute $\mathcal L_i^{(t)}$, $\mathcal H_i^{(t)}$, and the hidden states required by the anchor.
    \STATE Compute $\mathcal R_i^{(t)}$ by Eqs.~\eqref{eq:composite}--\eqref{eq:w}.
    \STATE Update $v_l^{(t)}$ and $\widetilde a_i^{(t)}$ by Eqs.~\eqref{eq:anchor-delta-mean}--\eqref{eq:anchor-projection-norm}.
    \STATE $s_i^{(t)}\gets G_i\mathcal R_i^{(t)}\bigl(1+\lambda\widetilde a_i^{(t)}\bigr)$.
    \STATE $\mathcal S_t\gets$ top-$B$ candidates by the key $\kappa_i^{(t)}$ of Eq.~\eqref{eq:tag-selection-key}.
    \STATE Fine-tune $\theta_{t-1}\to\theta_t$ on $\mathcal S_t$.
  \ENDFOR
  \STATE \textbf{return} $\theta_T$
  \end{algorithmic}
\end{algorithm}
